% Cover Letter: IXL Learning Software Engineer
% Author: Ajay Venkatesh

\documentclass[letterpaper,11pt]{article}

\usepackage[top=0.8in, bottom=0.8in, left=0.9in, right=0.9in]{geometry}
\usepackage[hidelinks]{hyperref}
\usepackage{lmodern}
\usepackage{parskip}
\usepackage{microtype}

\setlength{\parskip}{6pt}
\setlength{\parindent}{0pt}

\begin{document}

\begin{flushleft}
\textbf{\large Ajay Venkatesh} \\
Brooklyn, NY \\
+1 516 585 9013 \\
\href{mailto:av3855@nyu.edu}{av3855@nyu.edu}
\end{flushleft}

\vspace{10pt}

May 19, 2026

\vspace{6pt}

Hiring Manager \\
IXL Learning

\vspace{10pt}

Dear Hiring Manager,

My name is Ajay Venkatesh. I'm finishing my M.S. in Computer Engineering at NYU, with several years of production software experience: a few years at PwC in Bengaluru, a summer SDE internship at Amazon, and ongoing on-campus work at NYU's Novel AI Technologies. My focus has been on full-stack product engineering with strong \textbf{Java backends}, the data layer underneath, and shipping educational software that real students actually use.

The IXL posting is a strong match across a few axes. At Campus Mesh I built and continue to maintain an \textbf{EdTech product} that thousands of students depend on for academic workflows, including study groups, messaging, notifications, and profile management, backed by \textbf{Node.js microservices} and \textbf{MongoDB}. I also built \textbf{CampusIQ}, an LLM-powered AI assistant with \textbf{prompt engineering} and \textbf{RAG retrieval} over structured academic data, improving how students find the answers they need. At NYU's Novel AI Technologies I've spent the last year on on-campus product work with the same posture: long-running engagement with internal stakeholders, real users, and software that has to be correct because real people rely on it.

On the Java and backend side, at Amazon I engineered backend services in \textbf{Java} (Coral, Smithy, Dagger, CBOR) with low-level design patterns and unit/integration testing via \textbf{Mockito}, shipped under code-review discipline. I provisioned the supporting infrastructure with \textbf{AWS CDK} (Lambda, API Gateway, DynamoDB, IAM, CloudWatch) and operated multi-stage, multi-region CI/CD with rollback controls. At PwC I spent three years on a distributed backend on \textbf{Microsoft Azure} with event-driven processing and fault-tolerant retries over \textbf{SQL Server}, owning features from design through deployment alongside architects, PMs, and business stakeholders.

On full-stack, I've shipped \textbf{React} frontends at Amazon (React 18 with optimized API integration and client-side caching) and at NYU (analytics platform with role-based access for paying clients). I'm comfortable spanning back-end wiring, application logic, and UI in the same pull request, which is the shape of work the posting describes.

What pulls me toward IXL specifically is the mission. Education software runs on real consequences. Students whose learning depends on the system being correct, parents and teachers who trust the product to behave the way it claims. That's an environment where engineering judgment compounds, and where outcome-oriented thinking actually changes how learning happens.

I'd welcome the chance to discuss further. Thanks for considering my application.

\vspace{12pt}

Sincerely, \\
Ajay Venkatesh

\end{document}
